%-------------------------
% Resume — Ivan Vanichkin
% Python Backend Developer
%-------------------------

\documentclass[a4paper,10pt]{article}

\usepackage[utf8]{inputenc}
\usepackage[T2A]{fontenc}
\usepackage[russian]{babel}
\usepackage[empty]{fullpage}
\usepackage{titlesec}
\usepackage{enumitem}
\usepackage{fancyhdr}
\usepackage{tabularx}
\usepackage[hidelinks]{hyperref}

\pagestyle{fancy}
\fancyhf{}
\fancyfoot{}
\renewcommand{\headrulewidth}{0pt}
\renewcommand{\footrulewidth}{0pt}

\addtolength{\oddsidemargin}{-0.45in}
\addtolength{\evensidemargin}{-0.45in}
\addtolength{\textwidth}{0.9in}
\addtolength{\topmargin}{-0.65in}
\addtolength{\textheight}{1.3in}

\urlstyle{same}
\raggedbottom
\raggedright
\setlength{\tabcolsep}{0in}

\titleformat{\section}
  {\vspace{-5pt}\bfseries\raggedright\large}
  {}{0em}{}
  [\titlerule \vspace{-5pt}]

\newcommand{\resumeItem}[1]{
  \item\small{#1\vspace{-2pt}}
}

\newcommand{\resumeSubheading}[4]{
  \vspace{-1pt}\item
  \begin{tabular*}{\textwidth}{l@{\extracolsep{\fill}}r}
    \textbf{#1} & \textbf{#2} \\
    \textit{\small #3} & \textit{\small #4} \\
  \end{tabular*}\vspace{-6pt}
}

\newcommand{\resumeSubHeadingListStart}{\begin{itemize}[leftmargin=0pt,label={}]}
\newcommand{\resumeSubHeadingListEnd}{\end{itemize}}
\newcommand{\resumeItemListStart}{\begin{itemize}[leftmargin=17pt,itemsep=1pt,topsep=3pt]}
\newcommand{\resumeItemListEnd}{\end{itemize}\vspace{-5pt}}

\begin{document}

%----------HEADER----------

{\LARGE\textbf{Иван Ваничкин}}\\[3pt]
{\large\textbf{Junior Python Backend Developer}}\\[5pt]
\href{mailto:drdong654@gmail.com}{drdong654@gmail.com}
\enspace | \enspace +7~905~796~43~20
\enspace | \enspace \href{https://github.com/drdong654}{github.com/drdong654}
\enspace | \enspace \href{https://t.me/bdfy2003}{Telegram: @bdfy2003}

%----------SUMMARY----------

\section{О себе}

\small{
Junior Python Backend Developer с коммерческим опытом разработки backend-сервисов,
REST API и Telegram-приложений. Переквалифицировался в Python-разработку через
профильное обучение и практику, после чего начал работать над коммерческим продуктом.
Использую Python, FastAPI, PostgreSQL, SQLAlchemy, asyncio, Docker и pytest.
Имею практический опыт разработки RAG-систем, векторного поиска и интеграции LLM
в backend-приложения.
}

%----------EXPERIENCE----------

\section{Опыт работы}

\resumeSubHeadingListStart

\resumeSubheading
  {Vildly}
  {Python-разработчик}
  {Backend-разработка}
  {Апрель 2026 --- настоящее время}

\textbf{Стек:} Python, FastAPI, Aiogram 3, PostgreSQL, SQLAlchemy, asyncpg,
Docker, pytest, Git, RAG, pgvector

\resumeItemListStart
\resumeItem{Разрабатываю REST API и backend-логику с разделением API-слоя,
бизнес-логики и слоя доступа к данным.}
\resumeItem{Реализовал регистрацию и авторизацию пользователей, управление профилем,
валидацию входных данных и обработку ошибок.}
\resumeItem{Перевёл хранение пользовательских данных с файлового прототипа
на PostgreSQL и настроил асинхронный доступ через SQLAlchemy и asyncpg.}
\resumeItem{Разрабатываю Telegram-интерфейс на Aiogram 3 и FSM-сценарии,
интегрированные с backend-приложением.}
\resumeItem{Работаю над RAG-системой и векторным поиском по базе знаний
с использованием PostgreSQL и pgvector.}
\resumeItem{Пишу unit- и integration-тесты на pytest, контейнеризирую сервисы
с Docker и участвую в командной разработке через Git/GitHub.}
\resumeItemListEnd

\resumeSubheading
  {Управа района Текстильщики}
  {Ведущий специалист}
  {Автоматизация процессов; с апреля по июнь 2026 --- совмещение}
  {Июль 2025 --- Июнь 2026}

\resumeItemListStart
\resumeItem{Разрабатывал и внедрял Telegram-ботов для автоматизации отдельных
этапов рекрутмента и взаимодействия с участниками районных проектов.}
\resumeItem{Автоматизировал повторяющиеся операции, работал с корпоративными
информационными системами и структурированием рабочих данных.}
\resumeItem{Поддерживал контент сайта и контролировал актуальность
публикуемой информации.}
\resumeItem{Выполнил план по рекрутменту актива района на 100\%.}
\resumeItemListEnd

\resumeSubHeadingListEnd

%----------SKILLS----------

\section{Технические навыки}

\begin{itemize}[leftmargin=0pt,label={},itemsep=1pt,topsep=3pt]
\small{
\item \textbf{Backend:} Python, FastAPI, Aiogram 3, asyncio, REST API, ООП
\item \textbf{Базы данных:} PostgreSQL, SQL, SQLAlchemy, asyncpg, pgvector
\item \textbf{Тестирование:} pytest, pytest-asyncio, unit testing, integration testing
\item \textbf{AI / RAG:} RAG, LLM, Ollama, embeddings, vector search
\item \textbf{Инструменты:} Git, GitHub, Docker, Linux, Ruff, uv
}
\end{itemize}\vspace{-7pt}

%----------EDUCATION----------

\section{Образование}

\resumeSubHeadingListStart
\resumeSubheading
  {Российский новый университет}
  {2025 --- 2027}
  {Магистратура, Реклама и связи с общественностью}
  {}
\resumeSubheading
  {Российский новый университет}
  {2020 --- 2025}
  {Бакалавриат, Государственное и муниципальное управление}
  {}
\resumeSubHeadingListEnd

%----------RETRAINING----------

\section{Профессиональная переквалификация и обучение}

\resumeItemListStart
\resumeItem{\textbf{Нетология, 2026 --- «Основы Python: создаём телеграм-бота».}
Сертификат № PYFREE-ASYNC-10158846. Курс в рамках перехода в Python backend-разработку.}
\resumeItem{Практическая переквалификация продолжена в коммерческой backend-разработке:
FastAPI, PostgreSQL, асинхронный Python, тестирование, Docker и RAG.}
\resumeItemListEnd

%----------LANGUAGES----------

\vspace{2pt}
\textbf{Языки:} Русский --- родной, Английский --- B1

\end{document}
